\documentclass[twocolumn]{aastex631}

\usepackage{amsmath}
\usepackage{multirow}
\usepackage{natbib}
\usepackage{graphicx} 
\usepackage{aas_macros}

\begin{document}

Higher-derivative scalar-tensor theories, such as Degenerate Higher-Order Scalar-Tensor (DHOST) theories, provide a robust framework for modifying gravity, but their physical viability requires the absence of pathological ghost instabilities. A conjectured local, field-dependent symmetry is believed to be the underlying principle protecting the specific structure that ensures this classical health. This work directly investigates the fate of this protective symmetry in an effective action that includes quantum corrections, modeled by the addition of constant-coefficient Gauss-Bonnet and Weyl-squared terms to a ghost-free DHOST Lagrangian. We perform a rigorous, first-principles calculation of the variation of the complete action under the proposed symmetry transformation—a transformation that mixes the scalar and metric sectors—without any prior assumption of its invariance. Our analysis reveals a stark dichotomy: we confirm that the classical DHOST sector of the action is indeed invariant, provided the transformation parameters satisfy a specific differential condition, but we find that the curvature-squared terms explicitly break the symmetry. This result demonstrates a fundamental tension between the symmetry that guarantees classical stability and the standard form of leading-order quantum corrections, implying that a self-consistent, ghost-free quantum-corrected theory must incorporate such effects in a more subtle, symmetry-preserving fashion.

\end{document}
                